\scanpage{6}
\textit{C. R. Acad. Sc.}, t. 256, p. 842-844. \hfill Séance du 21 janvier 1963. \hfill Groupe 1

\textsc{Algèbre.} --- \textit{Sur les modules de différentielles.}

Note (*) de M. \textsc{Jean-Pierre Jouanolou}, présentée par M. René Garnier.

Application des modules de différentielles aux algèbres quasi-séparables et aux anneaux locaux réguliers.

$A$ et $B$ désignent dans toute cette Note deux anneaux commutatifs et unitaires, $u:A\to B$ un monomorphisme unitaire et $D_A(B)$ le $B$-module des différentielles de l'algèbre ainsi définie. Le passage de cette situation « globale » à la situation « locale » se fait en imposant à $A$, $B$, $u$ d'être locaux; l'idéal maximal de $A$ (resp. $B$) sera en ce cas noté $\mathfrak m$ (resp. $\mathfrak n$) et son corps résiduel $L_a$ (resp. $L_b$); par ailleurs, le rang de $\mathfrak m$ est $\operatorname{rg}\mathfrak m=\dim_{L_a}(\mathfrak m/\mathfrak m^2)$, et de même pour $\mathfrak n$. Rappelons qu'étant donnée une extension de corps : $K\to L$, la dimension du $L$-espace vectoriel $D_K(L)$ représente en caractéristique $0$ (resp. $p>0$) le degré de transcendance de $L$ par rapport à $K$ (resp. le nombre d'éléments d'une $p$-base de $L$ par rapport à $K$) : on la notera $\{L:K\}$ dans tous les cas.

Nous nous proposons d'étudier d'abord les cas où $D_A(B)$ est nul, puis ceux où il est libre (en liaison avec la régularité des anneaux locaux), le tout grâce aux deux suites exactes établies dans le cadre local par Berger et Kunz \textsuperscript{(1)} :

\diagram{\begin{tikzcd}
0\arrow[r]&\dfrac{\mathfrak n}{\mathfrak n^2+\mathfrak m B}\arrow[r]&L_b\otimes_B D_A(B)\arrow[r]&D_{L_a}(L_b)\arrow[r]&0,
\end{tikzcd}}

\diagram{\begin{tikzcd}
0\arrow[r]&\dfrac{\mathfrak n}{\mathfrak n^2+(\mathfrak n\cap A[B^p])B}\arrow[r]&L_b\otimes_B D_A(B)\arrow[r]&D_{L_a}(L_b)\arrow[r]&0.
\end{tikzcd}}

La première est vraie si $L_b$ est séparable sur $L_a$, la deuxième, sinon : alors $p>0$ désigne la caractéristique des corps résiduels.

1. \textsc{Structure des algèbres quasi-séparables.} --- Une $A$-algèbre $B$ est dite \uline{quasi-séparable} lorsque $D_A(B)=0$. Il s'agit là d'une propriété locale, en ce sens que $B$ est quasi-séparable si pour tout idéal maximal $\mathfrak n$ de $B$, $B_{\mathfrak n}$ est quasi-séparable sur $A_{\mathfrak m}$ ($\mathfrak m=\mathfrak n\cap A$), et seulement en ce cas.
\margin{\uline{\unsure{non ramifiée}}}

\textit{Localement :}

(i) Si $L_b$ est séparable sur $L_a$ (en particulier en caractéristique $0$) : lorsque $B$ est quasi-séparable,
\[
(\mathrm R)\quad \mathfrak n=\mathfrak n^2+\mathfrak mB,
\]
et $L_b$ est algébrique (resp. relativement parfait \textsuperscript{(2)}) sur $L_a$ en caractéristique $0$ (resp. $p>0$) la réciproque étant vraie si l'on suppose, en outre, que $D_A(B)$ est de type fini. Notons que si $B$ est noethérien, ou plus généralement si $\mathfrak n$ est de type fini, la relation (R) s'écrit
\[
\mathfrak n=\mathfrak mB,
\]
ce qui exprime que $B$ n'est pas ramifié sur $A$.
\margin{!?}

\mspage{5}
